\section{A taxonomy of surrogate models for energy system optimization: Model class, surrogate target, integration depth, and trust mechanism}
\label{sec:taxonomy}

FigureÂ~\ref{fig:surrogate-classification} illustrates the conceptual
role and classification scheme of surrogate models in energy system
optimization. It shows that their common application lies between
predictive modeling and optimization-based decision support. Surrogate
models can be classified by 1) model class, 2) surrogate target, 3)
optimizer role, 4) and trust mechanism. These four dimensions separate
the main questions that determine how a surrogate is used in an
optimization model: what mathematical approximation is employed, what
part of the original model is approximated, where the surrogate is
integrated in the optimization, and why its predictions can be trusted
for decision-making.

We propose this taxonomy because model class as the sole distinguishing
feature of surrogate models is not sufficient, because the same model
class can be employed for different purposes in an optimization model. A
Gaussian-process model may approximate an objective function, support
Bayesian infill sampling, quantify prediction uncertainty, or screen
candidate solutions before an expensive simulation is evaluated
\cite{Hu2020,Hu2021671,Volodina2022,Schulte2020}.
Similarly, a neural network may predict demand or renewable generation,
approximate a detailed building simulation or power-flow states, or
directly predict a near-optimal solution
\cite{Ruan2021221,Khaloie2025,Lim2025,Wu2022231,Chen2022,Chen20244723,Li202359995}.

\begin{table*}[!t]
\centering
\scriptsize
\begin{tabularx}{\textwidth}{@{}L{0.16\textwidth} C{0.05\textwidth} C{0.08\textwidth} L{0.21\textwidth} L{0.20\textwidth} L{0.22@\textwidth} Family & Output & Opt. compat. &@{}}
\toprule
Strengths & Typical failure modes & Where it dominates (refs.)
\\
Polynomial RSM / PCE & $p$,$i$ & $\nabla$ $\circ$ $\square$ &
Closed-form moments; analytic chance constraints; parsimonious & Curse
of dimensionality; smoothness assumption; sparse-grid cost &
Chance-constrained AC OPF, stochastic dispatch
\cite{Muhlpfordt20194806,Muhlpfordt20174490,Koirala20242284,Xu20212696,Engelmann20186188}
\\
Kriging / Gaussian process & $p$,$i$ & $\nabla$ & Predictive posterior;
Bayesian acquisition; small-sample use & $O(n^3)$ scaling;
stationarity assumption; kernel choice sensitive & Bayesian optimization
and uncertainty propagation
\cite{Nikolaidis2021,Hu2020,Hu2021671,Volodina2022,Schulte2020,Liu20222679}
\\
Radial basis functions & $p$ & $\nabla$ & Smooth interpolant; cheap;
metaheuristic-friendly & Centre placement is heuristic; weak in
extrapolation & Surrogate-assisted evolutionary OPF and HRES sizing
\cite{Liu20223726,Mahapatra2022721,Mahapatra20161921,Hussien20211883}
\\
Random forest / GBDT & $p$,$i$ & $\square$ & Mixed inputs; fast
classification and regression; quantile variants & Step responses break
gradients; piecewise-constant tails & Unit-status identification,
chance-constraint surrogates and thermal-interference screening
\cite{Liang20246508,Chen2023657,Petrova2025}
\\
Neural networks (MLP / CNN / RNN) & $p$ & $\nabla$ ($\square$) &
Universal approximation; large-data scaling; transfer learning &
Black-box; calibration drift; overconfidence on OOD & End-to-end OPF,
demand and price forecasting
\cite{Chen20244723,Chen2022,Wu2022231,Xu20223066,Guo20222057,Li202359995}
\\
Hard-constrained / structure-preserving NN & $p$ & $\nabla$ $\circ$
$\square$ & Encodes feasibility or structural constraints; can support
exact reformulation & Constraint design is problem-specific; training
and reformulation can be costly & Feasible dispatch proxies and
MILP-embedded conversion-loss surrogates
\cite{Chen20244723,Li202359995,Liu20242534}
\\
Graph neural networks & $p$ & $\nabla$ & Topology-aware; transferable
across grid sizes & Message-passing depth; oversmoothing on radial
grid trees & Wind-farm and multi-fidelity power-flow surrogates
\cite{Park2019,Taghizadeh2024,Khayambashi202453}
\\
Decision-focused / L2O proxies & $p$ & $\nabla$ & Trained for decision
quality, not regression error & Feasibility restoration needed; dataset
is optimizer-synthetic & Learning-to-optimize for SCED / SCUC / OPF
\cite{Wu2022231,Chen20244723,Chen2022,Li2023__2,Wu2024391,Liang20246508}
\\
\end{tabularx}
\end{table*}

\subsection{Model classes}
\label{sec:surrogate_classes}

TableÂ~\ref{tab:T1-taxonomy} compiles the most commong model classes in
the reviewed literature, the type of output they produce (point or
interval/posterior), their compatibility with the dominant optimizer
classes (gradient-based, convex, mixed-integer), and their typical
failure modes.

\subsubsection{Polynomial response surfaces and polynomial chaos expansions}
\label{sec:tax:rsm}

\paragraph{Polynomial response surface}

A polynomial response surface is a simple surrogate model that
approximates the relationship between design variables and a result
using a polynomial. Instead of running an expensive simulation or
optimization for each new combination of design variables, one can
estimate a function like:

$$\begin{equation}
 y \approx \beta_0 + \sum_{i=1}^{n} \beta_i x_i + \sum_{i=1}^{n} \beta_{ii} x_i^2 + \sum_{i<j}^{n} \beta_{ij} x_i x_j
\end{equation}$$

where $y$ denotes an output of interest, such as costs, CO~2~ emissions,
or efficiency, while $x_i$ represents the $i$-th design variable, such
as PV capacity, battery capacity, or electrolyzer power. The
coefficients $\beta$ are estimated from previously evaluated training
data.

Response-surface methodology (RSM) has been applied in PV control
parameter tuning \cite{Hasanien2018168}, thermal-system optimization
problems \cite{Sami20181416,Ahmad202451,Sun2020}, in chance-constrained
AC optimal power flow as an iterative response-surface refinement
\cite{Xu20212696}, and in green-hydrogen production models
\cite{Ahmad202451}.

Polynomial response surfaces have the advantage of clear
interpretability, ease of fit on small designs and direct compatibility
with gradient-based and convex solvers. A limitation is the curse of
dimensionality: the number of polynomial terms grows rapidly with the
number of inputs and polynomial order, while low-order response surfaces
may miss strong nonlinearities and interactions
\cite{Torun20192128,Hasanien2018168}.

\paragraph{Polynomial chaos expansion}

Polynomial chaos expansions (PCE) extend response surface methods to a
probabilistic approach by representing the model output as an expansion
in orthogonal polynomial basis functions of random input variables. For
example, considering the optimization of a PV--BESS--H$_2$ system, where
the output $Y$ denote annual system costs, CO~2~ emissions, or grid
imports, these outputs depend on uncertain input variables such as
electricity prices or load profiles. A polynomial chaos expansion
approximates the expensive optimization model as a function of these
uncertain inputs:

$$\begin{equation}
 Y \approx \sum_{\alpha \in \mathcal{A}} c_{\alpha} \Psi_{\alpha}(\boldsymbol{\xi})
\end{equation}$$ $$\begin{equation}
 \boldsymbol{\xi} = \left[ \xi_{\mathrm{price}}, \xi_{\mathrm{PV}}, \xi_{\mathrm{load}}, \xi_{\mathrm{H_2}}, \xi_{\mathrm{CAPEX}} \right].
\end{equation}$$

where, $Y$ denotes the model output, $\boldsymbol{\xi}$ is the vector of
random input variables, $\Psi_{\alpha}$ are multivariate orthogonal
polynomial basis functions, and $c_{\alpha}$ are expansion coefficients
estimated from model evaluations.

Mühlpfordt et al.Â~\cite{Muhlpfordt20194806} introduced the
chance-constrained AC optimal power flow using PCE and recent work
employed general PCE-based hosting-capacity calculations
\cite{Koirala20242284}, to data-driven distribution planning
\cite{Fu20202904,Fu2022}, and to uncertainty quantification in stochastic
economic dispatch \cite{Safta20172535}.

Because the expansion coefficients provide direct access to output
statistics, such as the mean and variance, without requiring additional
Monte Carlo sampling, PCE is particularly beneficial for uncertainty
propagation such as chance-constrained optimization AC optimal power
flow, hosting-capacity analysis and stochastic dispatch
\cite{Muhlpfordt20194806,Safta20172535,Xu20212696,Koirala20242284}.
Thus, PCE provide simple and interpretable approximations for design and
process optimization
\cite{Hasanien2018168,Sami20181416,Ahmad202451}.
However, the limitation is the rapid growth of the polynomial basis with
the number of stochastic input dimensions and by reduced accuracy for
non-smooth model responses
\cite{Muhlpfordt20194806,Safta20172535,Koirala20242284,Xu20212696}.

\subsubsection{Gaussian process regression and kriging}
\label{sec:tax:gp}

Gaussian process (GP) regression, also known as kriging in
geostatistics, is a widely established surrogate modeling approach for
optimization and Bayesian uncertainty quantification. Kriging originally
referred to optimal spatial interpolation based on a covariance
structure, whereas GP regression formulates the same idea in a
probabilistic machine-learning framework. In surrogate modeling, both
terms are commonly used for probabilistic surrogate models that provide
not only a point prediction, but also an associated predictive
uncertainty.

In an energy-system optimization problem, the input vector $\mathbf{x}$
may contain design variables such as PV capacity, battery capacity,
electrolyzer power, hydrogen storage capacity, and fuel-cell power:

$$\begin{equation}
 \mathbf{x} = \left[ P_{\mathrm{PV}}, E_{\mathrm{BESS}}, P_{\mathrm{ELY}}, E_{\mathrm{H_2}}, P_{\mathrm{FC}} \right].
\end{equation}$$

The corresponding model output $y$ may denote annual system costs, CO~2~
emissions, grid imports, or another performance indicator obtained from
a detailed optimization model. A key advantage of GP regression is that
it predicts both the expected system performance and the uncertainty of
this prediction:

$$\begin{equation}
 \hat{y}(\mathbf{x}) = \mu(\mathbf{x}) \pm \sigma(\mathbf{x}).
\end{equation}$$

Here, $\hat{y}(\mathbf{x})$ denotes the surrogate prediction for a given
energy-system design, $\mu(\mathbf{x})$ is the predicted mean
performance, and $\sigma(\mathbf{x})$ represents the predictive
uncertainty. In practical terms, the mean prediction indicates how
promising a design is, while the uncertainty estimate indicates how
confident the surrogate is in this prediction.

In dispatch and operations, Nikolaidis et al.Â~\cite{Nikolaidis2021} used
GP-based Bayesian optimization, demonstrating substantial reductions in
the number of MILP evaluations. Hu et al.Â~\cite{Hu2020,Hu2021671} use GP
regression to propagate uncertainty in stochastic economic dispatch,
with a Bayesian-network extension in \cite{Volodina2022}. Wang et
al.Â~\cite{Wang2023} employ GP regression in an LSTM-based hybrid for
wind-power forecasting and dispatch, similar to that used by Zhang et
al.Â~\cite{Zhang2019270}. In Bayesian power-network planning \cite{Lawson201647}
GPs often serve as a surrogate model for transmission expansion under
uncertainty. In design and capacity expansion optimization studies,
kriging-based surrogate modeling has been used to solve large optimal
power flows \cite{Deng2020831} and to optimize battery thermal management
with GP-based screening \cite{Li2021,Babaeiyazdi2021,Kong2021}.
GP and kriging models are also used in surrogate-assisted design search
and calibration, including geothermal design, building and HVAC
calibration, and interactive multi-objective building-energy design
\cite{Schulte2020,Chakrabarty2021,Aghaei Pour2022303}.

A major limitation of GP regression is its computational cost. Standard
GP training becomes increasingly expensive as the number of evaluated
simulation points grows, because the model has to account for the
correlation between all training data points, making it difficult to
apply to very large training data sets. This problem can be addressed by
reducing the amount of data or correlation structure that the GP has to
process. For example, sparse kernel methods simplify the covariance
structure by assuming that only the most relevant or nearby training
points have a strong influence on a prediction. Moreover, local GP
models can be used to split the design space into smaller regions and
fit separate Gaussian processes within these regions, so that each model
only uses a subset of the full training data. Inducing-point methods
replace the full training set with a smaller set of representative
support points, which summarize the main structure of the data. These
approaches make GP regression more scalable when the number of runs
grows, but they are also accompanied by a loss of approximation
accuracy \cite{Hu2020,Volodina2022,Liu20222679}.

\subsubsection{Radial-basis-function networks and kernel regressors}
\label{sec:tax:rbf}

Radial-basis-function (RBF) models are deterministic surrogate models
that combine several localized basis functions (functions that have a
strong effect in a specific area of the input/design space and less in
another). The basic idea is that a new design is predicted from its
similarity to previously selected support points. Designs close to a
support point have a strong influence on the prediction, while distant
designs have only a weak influence. An RBF surrogate approximates this
output as a weighted sum of radial basis functions:

$$\begin{equation}
 \hat{y}(\mathbf{x}) = \sum_{j=1}^{M} w_j \, \phi\left(\left\| \mathbf{x} - \mathbf{c}_j \right\|\right).
\end{equation}$$

Here, $\mathbf{c}_j$ denotes the $j$-th support point, $w_j$ is its
fitted weight, and $\phi(\cdot)$ is a radial basis function that depends
on the distance between the new design $\mathbf{x}$ and the support
point $\mathbf{c}_j$.

A common choice is the Gaussian radial basis function:

$$\begin{equation}
 \phi\left(\left\| \mathbf{x} - \mathbf{c}_j \right\|\right) = \exp\left( -\frac{\left\| \mathbf{x} - \mathbf{c}_j \right\|^2}{2\sigma^2} \right).
\end{equation}$$

This means that nearby energy-system designs have a stronger effect on
the prediction than distant designs. Compared with Gaussian process
regression, RBF models use a similar distance-based function, but they
usually provide a deterministic prediction rather than a full posterior
uncertainty distribution. This makes them computationally less expensive
when the surrogate has to be evaluated many times, for example during
candidate sampling.

Urgun et al.Â~\cite{Urgun20202791} use multilabel RBF classification for
importance-sampled composite-power-system reliability, and
Karamichailidou et al.Â~\cite{Karamichailidou20212137} use RBF to model
wind-turbine power curves for economic dispatch. RBF is also used in PV
production \cite{Yang2020415,Tian2020}, and have been used as the
regression back-end in extreme-learning-machine load forecasting
\cite{Zeng2017175,Sideratos2020,Madhukumar20228891,Rafi202132436}. In
building-energy applications, SVR and ANN metamodels have been used to
predict district cooling loads \cite{Jia2022553}, while an ANN trained on
EnergyPlus simulations has supplied fitness values to a genetic
algorithm for heating set-point scheduling \cite{Reynolds2017704}.

RBF models show fast evaluation and work well with relatively small
training data sets. However, their predictions are only reliable within
the region covered by the evaluated simulation points. For new
energy-system designs that are far away from the training data, e.g.
very different combinations of PV capacity, battery capacity,
electrolyzer power, or hydrogen storage size, RBF models show rather
poor approximation. Therefore, RBF-based optimization is often combined
with active-sampling strategies to add new simulations in promising or
poorly covered parts of the design space.
\cite{Deng2020831,Tian2020,Han2018622,Pang2023}.

\subsubsection{Tree ensembles: random forests and gradient boosting}
\label{sec:tax:trees}

Tree ensembles combine decision trees to obtain more accurate
predictions than a single tree. They are widely used for regression and
classification tasks in energy-system optimization models for different
tasks because due to their ability to perform well on structured tabular
data and can capture nonlinear relationships and interactions between
input variables. In tree-ensembles, the output is a continuous
performance indicator, such as costs or CO~2~ emissions. In
classification tasks, the output is a discrete class label, for example
whether an energy-system design is feasible or infeasible.

For a random forest, the prediction is obtained by averaging the outputs
of many individual decision trees:

$$\begin{equation}
 \hat{y}(\mathbf{x}) = \frac{1}{B} \sum_{b=1}^{B} T_b(\mathbf{x}),
\end{equation}$$

where $T_b(\mathbf{x})$ is the prediction of the $b$-th tree and $B$ is
the number of trees in the ensemble. In classification tasks, such as
predicting whether an operating point infeasible, the individual trees
vote for a class instead of returning a continuous value.

Gradient-boosted trees use a different principle: Instead of training
all trees independently, they build the ensemble sequentially where each
new tree is trained to correct the remaining error of the previous
trees:

$$\begin{equation}
 \hat{y}(\mathbf{x}) = \sum_{m=1}^{M} \eta_m T_m(\mathbf{x}),
\end{equation}$$

where $T_m(\mathbf{x})$ denotes the $m$-th tree and $\eta_m$ is its
learning weight. XGBoost is a widely used implementation of this
gradient-boosting principle. Tree-based ensemble methods have been used
in microgrid scheduling to serve as fast predictive models that
approximate the performance or feasibility of candidate operating
schedules, thereby reducing the number of detailed simulation or
optimization runs required,
\cite{Eseye201991463,Faraji2020157284,Wang20192635,Nematollahi2021},
distribution-system planning under uncertainty
\cite{Fu20202904,Fu2022,Jahangiri20244849,Jahangiri2023}, and
short-term wind, solar and load forecasting that determines optimization
dispatch \cite{Li2020,Hanifi2024,Alkesaiberi2022,Madhukumar20228891}.
Also, tree ensembles have been used in surrogate-assisted MOO where they
classify candidate decisions as feasible or infeasible before more
expensive optimization
\cite{Han2018622,Schulte2020,Pang2023,Xiao2018}.

The main advantage of tree ensembles is that they can model complex
nonlinear effects without requiring an explicit functional form. For
example, they can learn how the effect of battery capacity depends on PV
capacity, electricity prices, or hydrogen storage size. A limitation of
tree-based models is, however, that their predictions are
piecewise-constant or piecewise-linear and therefore do not provide
smooth gradients. This means that they change their prediction abruptly
at decision boundaries instead of forming a smooth function; therefore,
continuous derivatives, as required by gradient-based optimizers, cannot
be used. This makes them difficult to embed directly into gradient-based
nonlinear programming problems. In practice, tree models are therefore
often used outside the optimization model as fast prediction models, for
example to evaluate candidate solutions before or after
optimization.\cite{Ren2024,Lin2023,Pang2023,Wang20211767}.

\subsubsection{Neural surrogates: MLP, recurrent and graph architectures}
\label{sec:tax:nn}

Neural networks are widely used as surrogate models in recent
energy-system literature because they can approximate complex nonlinear
relationships between inputs and outputs. In contrast to simpler
surrogate models, neural networks can learn interactions between many
design variables, operating conditions, and time-dependent system
states. This makes them especially beneficial for highly complex
nonlinear optimization models.

A neural network-based surrogate learns an approximate relationship from
inputs to outputs:

$$\begin{equation}
 \hat{y} = f_{\theta}(\mathbf{x}),
\end{equation}$$

where $f_{\theta}$ denotes the neural network and $\theta$ represents
its trained weights.

For static design problems, multilayer perceptrons (MLPs) are commonly
used. They relate a fixed input vector, such as technology capacities
and cost assumptions, to one or more performance indicators:

$$\begin{equation}
 \hat{y} = f_{\theta} \left( P_{\mathrm{PV}}, E_{\mathrm{BESS}}, P_{\mathrm{ELY}}, E_{\mathrm{H_2}}, P_{\mathrm{FC}} \right).
\end{equation}$$

For time-dependent problems, recurrent neural networks can be used to
learn from input trajectories, such as hourly demand, renewable
generation, prices, or storage states:

$$\begin{equation}
 \hat{y}_t = f_{\theta} \left( \mathbf{x}_1,\mathbf{x}_2,\ldots,\mathbf{x}_t \right).
\end{equation}$$

This is useful when the surrogate has to approximate dispatch behaviour,
storage operation, or system states over time. Graph neural networks are
particularly relevant for networked energy systems because they can
represent the physical structure of electricity, gas, or heat networks.
In this case, buses, substations, buildings, or heat network nodes are
represented as graph nodes, while power lines, pipes, or thermal
connections are represented as edges. A graph-based surrogate can
therefore learn how local injections, loads, or temperatures evolve
throughout the network:

$$\begin{equation}
 \hat{\mathbf{y}} = f_{\theta} \left( \mathbf{X}, \mathbf{A} \right),
\end{equation}$$

where $\mathbf{X}$ contains node-level features such as demand,
generation, temperature, or pressure, and $\mathbf{A}$ describes the
network topology.

MLPs are used when the surrogate receives a fixed-size input vector and
returns a fixed-size output vector. This is typical for optimization
problems in which variables such as technology capacities, demand
levels, or prices are used to predict performance indicators (e.g. costs
or emissions). Because trained MLPs are fast to evaluate, they are
useful when many candidate solutions have to be assessed during an
optimization.

In the reviewed literature, MLPs have been used as fast approximations
of costs, feasibility, or security constraints in security-constrained
economic dispatch
\cite{Chen2022,Chen20244723,Wu2022231,Xu20223066,Guo20222057,Li202359995}.
They are also used in response-surface models in metamodel-based design
search, where they replace repeated evaluations of more detailed
optimization models
\cite{Hirvonen2017323,Jia2022553,Li2023,Wang2020202933}.
In the building sector, MLPs are used as surrogates for retrofit
analysis, for example to approximate energy demand, costs, or other
building-performance indicators
\cite{Thrampoulidis2021,Reynolds2017704,Koschwitz2018134}. Recurrent and
convolutional neural networks are used when the input is sequential
rather than static, meaning that the surrogate has to learn from ordered
time-series data, such as hourly demand, renewable generation,
electricity prices, or storage states. This includes long short-term
memory (LSTM) networks, gated recurrent units (GRUs), convolutional
neural networks (CNNs), and hybrid convolutional neural network--long
short-term memory (CNN--LSTM) models. They are widely used in dispatch
and storage scheduling, where current decisions depend on previous
system states and operating conditions
\cite{Rafi202132436,Sideratos2020,Li2019,Tian2020,Yang2020415,Wang2023,Shi2024,Lu20232989,Hanifi2024}.
Graph neural networks (GNNs) are increasingly used for network-based
energy-system optimization problems. They are well suited to analyze
power flows, voltage magnitudes because they can use the topology of the
underlying grid directly. In this representation, buses, substations, or
other network nodes are treated as graph nodes, while lines or pipes are
represented as graph edges
\cite{Park2019,Su2022315,Liu20223726,Chen2023657}. Where the
underlying physical equations are well understood, physics-informed
neural networks (PINNs) can include these equations directly in the
training data. This can reduce extrapolation error, because the
surrogate is not only fitted to data but also encouraged to respect
known physical constraints
\cite{Park2019,Su2022315,Liu20222679,Chen2023657}.

Overall, neural surrogates can represent nonlinear, temporal, and
especially network-dependent system behaviour. Their main limitations
are the need for sufficient training data, weaker interpretability
compared with simpler models, and the risk of poor generalization
outside the range of the training data.

\subsection{Surrogate targets}
\label{sec:tax:targets}

Surrogate models can also be classified by their target, i.e. what the
surrogate approximates. Some surrogates approximate exogenous inputs
such as load, weather, renewable energy generation or prices, which
serve as an input to the optimization model
\cite{Eseye201991463,Li2020,Hanifi2024,Conti2026}. Objective values
can be approximated directly, such as energy demand, cost, or emissions
\cite{Westermann2019,Hirvonen2017323,Ahmad202451}, or
physical states or constraints, such as voltage magnitudes, line flows,
storage states, indoor temperatures, thermal-network states or
feasibility violations can be emulated
\cite{Park2019,Su2022315,Mohammadi2024,Khaloie2025}. Surrogates can
also approximate probability distributions, posterior variance,
quantiles or chance-constraint terms
\cite{Muhlpfordt20194806,Hu2020,Volodina2022,Tan2026}, and
learning-to-optimize approaches approximate decisions directly, for
example dispatch vectors, active sets, optimal power flow solutions, or
dispatch schedules
\cite{Wu2022231,Chen2022,Chen20244723,Xu20223066,Li202359995,Liang20246508}.

\subsection{Optimizer integration depth}
\label{sec:tax:integration-depth}

Surrogate models can also be classified according to integration depth
in the optimization.
FigureÂ~\ref{fig:taxonomy-integration-depth} shows how the reviewed
literature is distributed across optimizer integration depth and how it
relates to surrogate target, model class and trust mechanism. In a weak
integration, the surrogate provides forecasts, physical data, or
preprocessed inputs outside of the optimization model itself
\cite{Eseye201991463,Li2020,Hanifi2024}. Surrogates can also preselect
candidate solutions or identify promising regions before expensive model
evaluations are used \cite{Han2018622,Schulte2020,Pang2023,Xiao2018}
or replace a computationally expensive objective function or physical
simulations, and perform constraint checks
\cite{Hirvonen2017323,Deng2020831,Thrampoulidis2021}. In
Bayesian optimization and surrogate-assisted evolutionary algorithm
optimization, surrogate models can balance the exploration of promising
and uncertain search spaces
\cite{DiazManriquez2016,Hu2020,Schulte2020,Volodina2022}. Surrogate
can also be embedded into the optimization problem, for example as a
mixed-integer-compatible approximation
\cite{Wu2022231,Chen2022,Lin2023,Pang2023}. At the strongest
integration depth, learning-to-optimize models replace the optimizer by
predicting decisions directly
\cite{Wu2022231,Chen20244723,Li202359995,Liang20246508}.

\subsection{Trust mechanisms and optimizer-compatible surrogates}
\label{sec:tax:trust-mechanisms}

Another taxonomy dimension is the trust mechanism. Usually, in black-box
surrogate models, which are purely data-driven approximations whose
internal structure is not constrained by physical equations or
optimization constraints, prediction errors on validation data, mostly
Root Mean Squared Error (RMSE), Mean Absolute Error (MAE) or the
coefficient of determination $R^2$ are used to validate the surrogate
model \cite{Westermann2019,Zhang2026_building,Elwy2024}. For many
optimization studies, however, these predictive accuracy metrics alone
are often insufficient: a small average error can still lead to wrong
rankings of candidate solutions, infeasible solutions, or constraint
violations \cite{DiazManriquez2016,Ruan2021221,Khaloie2025,Lim2025}.

Some surrogate classes, therefore, provide additional trust mechanisms.
Gaussian processes and Bayesian models provide posterior uncertainty
that can be used for acquisition functions, which determine the next
optimal evaluation point based on model predictions
\cite{Hu2020,Hu2021671,Tan2026}. Polynomial chaos expansions provide
uncertainty propagation for probabilistic constraints
\cite{Muhlpfordt20194806,Safta20172535,Koirala20242284}, and
physics-informed neural network (PINN) models are validated through
balance equations, physical structure or network topology
\cite{Park2019,Su2022315,Mohammadi2024,Petrova2025,Du2025}. Thus,
decision-focused surrogates should be assessed by their effect on the
optimization results, such as feasibility rate, Pareto-front quality,
and constraint violation in relationship to computational acceleration
\cite{Ruan2021221,Khaloie2025,Lim2025,Chen20244723}.

\subsubsection{Optimization-aware surrogate formulations}
\label{sec:tax:physics-guided}

The preceding subsection classified surrogates mainly by model class.
For optimization, however, the model class alone is often not the most
important property. A surrogate with low prediction error can still lead
to infeasible solutions or distorted Pareto fronts when it is used
inside an optimizer
\cite{DiazManriquez2016,Ruan2021221,Khaloie2025,Lim2025}. Therefore,
some surrogate approaches are better described by how they make the
approximation reliable enough for optimization. In the reviewed
literature, three such optimization-aware surrogate formulations are
particularly relevant: physics-guided modeling, structure-preserving
formulations and decision-oriented surrogates.

\paragraph{Physics-guided and hybrid surrogates}

Physics-guided and hybrid surrogates use available physical knowledge
and information instead of learning the full input-output relation from
data alone. This knowledge may enter the optimization model through
simplified physical equations. The purpose is to improve computational
efficiency and extrapolation, especially when training data are sparse
or when the optimizer evaluates the surrogate outside the original
training range.

A common form is residual learning:

$$\begin{equation}
 \hat{y}(x) = f_{\mathrm{phys}}(x) + r_{\theta}(x),
\end{equation}$$

where $f_{\mathrm{phys}}$ is a simplified physical model and
$r_{\theta}$ is a learned correction term. The physical model captures
the main system behaviour, while the residual model learns the remaining
mismatch. This approach has been used in power-system, multi-energy and
sector-coupled applications where a complete high-fidelity model is too
expensive for repeated use inside the optimizer
\cite{Park2019,Liu20223726,Chen2023657,Liu20222679,Ren2024,Jalving2023,Zheng20231907}.
Physics-informed neural networks are one specific case of this broader
class. Their defining feature is not only the neural-network
architecture, but the inclusion of physical residuals in the training
loss. In energy-system models, such residuals may represent power
balance, heat balance, storage dynamics, network-flow equations or
conversion constraints. Related structure-preserving approaches embed
physical balances or known network structure directly into the
surrogate architecture or optimization layer
\cite{Park2019,Liu20223726,Chen2023657,Liu20242534}.
The main advantage is improved physical plausibility; the main
limitation is that the simplified physics and the residual weights must
be chosen carefully, otherwise the surrogate may become biased or
difficult to train.

\paragraph{Structure-preserving and solver-compatible surrogates}
\label{sec:tax:structured}

Structure-preserving and solver-compatible surrogates are designed so
that the approximation retains properties needed by the optimization
model. These properties may include convexity, monotonicity, feasibility
preservation, active-set consistency or compatibility with mixed-integer
linear formulations. This is especially important when the surrogate
approximates constraints, feasibility regions, cost curves, conversion
losses or operational decisions, because small prediction errors can
directly change the optimal solution. Convex or monotone surrogates are
designed so that they do not only fit the training data, but also
preserve known shape properties of the original model. For example, if
fuel use always increases when output power increases, the surrogate
should also be monotone. If operating costs increase in a convex way
with load or capacity, the surrogate should preserve this convex shape.
Such restrictions make the surrogate less flexible than a fully
black-box model, but they reduce the risk that the optimizer exploits
physically implausible artefacts of the approximation
\cite{Wang2020202933,Hirvonen2017323,Liu20223726,Chen2023657,Chen20244723}.

Other approaches aim at direct solver compatibility. For example,
piecewise-linear neural networks, especially ReLU-based networks (i.e.,
networks using the activation function $\mathrm{ReLU}(z)=\max(0,z)$,
which makes the learned mapping piecewise linear), can in some cases be
reformulated as mixed-integer linear constraints and embedded into a
MILP. This can make the surrogate usable inside an optimization model,
but also increase the model size through additional binary variables and
big-$M$ constraints, i.e. constraints that use a sufficiently large
constant $M$ to switch linear constraints on or off depending on binary
variables.\cite{Wu2022231,Chen2022,Xu20223066,Lin2023}.

A related group of methods enforces hard linear constraints, predicts
active sets, or generates near-feasible solutions that the optimizer can
validate, repair or refine \cite{Chen20244723,Li202359995}. These
surrogates are therefore not only judged by prediction accuracy. They
are useful because they preserve enough of the original optimization
model to avoid poor prediction accuracy. Their main limitation is the
trade-off between structure and flexibility: stronger restrictions can
improve feasibility and solver compatibility, but may reduce
approximation accuracy or make the embedded optimization problem larger.

\subsubsection{Decision-oriented surrogates and learning-to-optimize}
\label{sec:tax:l2o}

Instead of approximating an intermediate model output, decision-oriented
surrogates approximate the solution of an optimization problem, or at
least a solution-related object such as a dispatch vector. In this case,
the surrogate target is no longer only a system response or constraint
term, but the decision itself.

For example, an optimization problem relates data $p$, such as demand,
renewable generation, network state or prices, to an optimal solution:

$$\begin{equation}
 x^{\star}(p) = \arg\min_{x \in \mathcal{X}(p)} f(x,p).
\end{equation}$$

A learning-to-optimize surrogate learns an approximation of this
relation

$$\begin{equation}
 \hat{x}_{\theta}(p) \approx x^{\star}(p),
\end{equation}$$

and can then be used as a direct solution proxy.

This approach is particularly prominent in power-system optimization. Wu
et al.Â~\cite{Wu2022231} train a deep neural network to solve
security-constrained unit commitment with uncertain wind power. Chen et
al.Â~\cite{Chen2022,Chen20244723} learn optimization proxies for
large-scale security-constrained economic dispatch with feasibility
guarantees. Guo et al.Â~\cite{Guo20222057} learn optimal power allocation in
islanded microgrids, and Xu et al.Â~\cite{Xu20223066} train an ensemble deep
network for non-convex economic dispatch. Li et al.Â~\cite{Li202359995} learn
to solve optimization problems with hard linear constraints, while Liang
et al.Â~\cite{Liang20246508} combine learning-to-optimize with statistically
feasible optimization for chance-constrained dispatch.

Recent reviews show the rapid growth of this in power-system modeling
\cite{Khaloie2025,Lim2025,Ruan2021221}. If the surrogate predicts a
feasible or near-feasible solution, the model can avoid solving a large
MILP or dispatch problem from scratch, or can at least start the solver
from a much better initial point. The main risk is that the surrogate
now directly affects the decision. Its validation must therefore go
beyond the basic prediction error metrics such as RSME, MAE, and $R^2$,
and include feasibility rate, constraint violation, and robustness using
sensitivity analysis.
